\documentclass[11pt,a4paper]{article}
% Common preamble for the hint sheets (\input from each file)
\usepackage[utf8]{inputenc}
\usepackage[T1]{fontenc}
\usepackage{lmodern}
\usepackage{amsmath,amssymb,amsthm,mathtools}
\usepackage[a4paper,margin=2.3cm]{geometry}
\usepackage{enumitem}
\usepackage{xcolor}
\usepackage[hidelinks]{hyperref}
\renewcommand{\leq}{\leqslant}
\renewcommand{\le}{\leqslant}
\renewcommand{\geq}{\geqslant}
\renewcommand{\ge}{\geqslant}
\setlength{\parindent}{0pt}
\setlength{\parskip}{0.35em}
% \hintstitle{Sheet title}{Class work / Homework}
\newcommand{\hintstitle}[2]{\begin{center}{\LARGE\bfseries Hints: #1\par}\vspace{0.5em}{\large #2\par}\end{center}
\noindent\textit{The hints to each problem are cumulative. Read one, try again, and open the next only if you are still stuck.}\par\medskip\hrule\medskip}
\newcommand{\problem}[1]{\par\bigskip{\large\bfseries Problem #1}\par\nopagebreak\smallskip}
\newcommand{\subpart}[1]{\par\smallskip\textbf{(#1)}\par\nopagebreak}
\newcommand{\hint}[1]{\par\smallskip\noindent\textcolor{black!70}{\textbf{Hint #1.}}\enspace}
\newcommand{\note}[1]{\par\smallskip\noindent\textit{Note.} #1\par}
\providecommand{\lt}{<}
\providecommand{\gt}{>}
\providecommand{\tightlist}{\setlength{\itemsep}{0pt}\setlength{\parskip}{0pt}}
\emergencystretch=3em

\begin{document}
\hintstitle{Number theory -- 5}{Homework}

\problem{1}
\hint{1} For $x^2+y^2+z^2=2xyz$, look modulo $4$. The right side is even, so an even number of $x,y,z$ are odd. If exactly two are odd, the left side is $\equiv2\pmod4$ while the right side is $\equiv0$. Hence all three are even.
\hint{2} Writing $x=2x'$ and so on gives $x'^2+y'^2+z'^2=4x'y'z'$. Prove more generally that $x^2+y^2+z^2=2^kxyz$ ($k\ge1$) forces $x,y,z$ even and leads to the same equation with $k+1$. Infinite descent finishes the proof.
\hint{3} Markov: fix $x,y$. The two roots of $t^2-3xyt+x^2+y^2=0$ are $z$ and $z'=3xy-z=(x^2+y^2)/z$, so $z'$ is a positive integer. For $x\le y\le z$ with $(x,y)\ne(1,1)$, the quadratic is negative at $t=y$, so $z'<y\le z$.
\hint{4} So every solution descends by such jumps to one with $x=y=1$, i.e.\ to $(1,1,1)$ or $(1,1,2)$, and hence to $(1,1,1)$. The solutions are exactly the triples obtained from $(1,1,1)$ by repeatedly replacing one coordinate $z$ by $3xy-z$, together with their permutations (the Markov tree).

\problem{2}
\hint{1} View the equation as a quadratic in $m$. Its two roots $m,m'$ satisfy $mm'=2\cdot3^n$ and $m+m'=2^{n+1}-1$. Since the sum is odd, the roots are $3^p$ and $2\cdot3^q$ with $p+q=n$.
\hint{2} Let $k=\min(p,q)$. Then $3^k\mid2^{n+1}-1$, which forces $2\cdot3^{k-1}\mid n+1$ when $k\ge1$ (the order of $2$ mod $3^k$). On the other hand the larger root is $\ge3^{n-k}$ and is less than $2^{n+1}$, so $3^{n-k}<2^{n+1}$.
\hint{3} Combining $k>n-(n+1)\log_32$ with $2\cdot3^{k-1}\le n+1$ leaves only $n\le5$. Check these by hand. Answer: $(m,n)=(6,3),(9,3),(9,5),(54,5)$.

\problem{3}
\hint{1} Necessity: multiply the second equation by $3^M$, $M=\max a_k$. You get $\sum k\,3^{M-a_k}=3^M$. Reduce mod $2$: $1+2+\dots+n$ must be odd, so $n\equiv1,2\pmod4$.
\hint{2} For sufficiency, build solutions by splitting terms. A term $k$ with exponent $a$ can be replaced by the terms $k$ and $2k$, both with exponent $a+1$. Check that both sums are unchanged. This adds a new \emph{even} index $2k$; in particular $n=4k+1$ gives $n+1$.
\hint{3} To add odd indices, use a three-term split: replace term $m$ (exponent $a$) by $m$ with exponent $a+2$, $2m-1$ with exponent $a+1$, and $2m+3$ with exponent $a+2$. Check $\frac14+\frac12+\frac14=1$ and $\frac m9+\frac{2m-1}3+\frac{2m+3}9=m$.
\hint{4} Go from $n=4k+1$ to $n+8$. The odd indices $4k+3,4k+7$ come from $m=2k+2$, and $4k+5,4k+9$ from $m=2k+3$. The even indices $4k+2,\dots,4k+8$ come from doubling (do the splits in a suitable order). Base cases: $n=1$: $(0)$; $n=5$: $(a_1,\dots,a_5)=(2,1,3,4,4)$. Then $n\equiv2$ follows by one doubling.

\problem{4}
\hint{1} Induct on $k$; for $k=1$ take $m_1=n$. The idea is to write $1+\frac{2^k-1}n$ as one factor $1+\frac1m$ times an expression of the same form with $k-1$.
\hint{2} If $n=2t-1$ is odd: $1+\frac{2^k-1}{n}=\Bigl(1+\frac{2^{k-1}-1}{t}\Bigr)\Bigl(1+\frac1{n}\Bigr)$.
\hint{3} If $n=2t$ is even: $1+\frac{2^k-1}{n}=\Bigl(1+\frac{2^{k-1}-1}{t}\Bigr)\Bigl(1+\frac1{2t+2^k-2}\Bigr)$. Verify both identities and apply the induction hypothesis to the first factor.

\problem{5}
\hint{1} Compute small cases ($n=2$: $\{2,3\}$, $n=3$: $\{4,6,7\}$) and guess. Answer: $(n-2)2^n+1$.
\hint{2} Representability of every larger $m$, by induction on $n$: if $m$ is even, represent $m/2$ by $A_{n-1}$ and double; note that $2(2^{n-1}-2^k)\in A_n$. If $m$ is odd, do the same with $m-(2^n-1)$. Check that the halves exceed $(n-3)2^{n-1}+1$.
\hint{3} Non-representability of $N=(n-2)2^n+1$: if $N$ is a sum of $s$ elements $2^n-2^{k_i}$ ($0\le k_i\le n-1$), then $\sum2^{k_i}=s2^n-N=t\,2^n-1$ with $t=s-n+2\ge1$. Show that writing $t2^n-1$ as a sum of powers of $2$ not exceeding $2^{n-1}$ needs at least $(2t-1)+(n-1)$ terms. That is more than $s=t+n-2$, a contradiction.

\problem{6}
\hint{1} If $u=x+y\sqrt2>1$ has norm $\pm1$, then $u(\sqrt2-1)=(2y-x)+(x-y)\sqrt2$ is a smaller unit, still $\ge1$. Descend to $1$. All solutions of $x^2-2y^2=\pm1$ are $\pm(1+\sqrt2)^k$, $k\in\mathbb Z$; the norm is $(-1)^k$.
\hint{2} $\frac{m(m+1)}2=s^2$ is equivalent to $(2m+1)^2-2(2s)^2=1$. So $(2m+1)+2s\sqrt2=(3+2\sqrt2)^j$. This gives $s=1,6,35,204,\dots$, i.e.\ the square triangular numbers $1,36,1225,41616,\dots$
\hint{3} General $d$: use $p_{k-1}^2-dq_{k-1}^2=(-1)^kQ_k$ with $Q_k=1$ exactly when the period $r$ divides $k$. If $r$ is odd, then $k=r$ gives norm $-1$.
\hint{4} Conversely, any positive solution of $x^2-dy^2=-1$ has $x/y$ a convergent (Legendre, since $1<\sqrt d$). So $-1=(-1)^kQ_k$ with $Q_k=1$, hence $r\mid k$ and $k$ is odd. Therefore $r$ is odd.

\problem{7}
\hint{1} If $x$ is even, then $y^2\equiv7\pmod8$, which is impossible. If $x\equiv3\pmod4$, then $y^2\equiv2\pmod4$, also impossible. So $x\equiv1\pmod4$.
\hint{2} Write $y^2+1=x^3+8=(x+2)(x^2-2x+4)$. Here $x^2-2x+4=(x-1)^2+3\equiv3\pmod4$, so it has a prime factor $p\equiv3\pmod4$. Then $p\mid y^2+1$, which is impossible.
\hint{3} For $y^2=x^3-2$: $y$ is odd, and in $\mathbb Z[\sqrt{-2}]$ (a UFD with units $\pm1$) $(y+\sqrt{-2})(y-\sqrt{-2})=x^3$ with coprime factors. A common prime factor would divide $2\sqrt{-2}$ and hence $2$, but $x$ is odd. So $y+\sqrt{-2}=(a+b\sqrt{-2})^3$.
\hint{4} Comparing the $\sqrt{-2}$-parts gives $1=b(3a^2-2b^2)$, so $b=1$ and $a=\pm1$. Then $y=a^3-6ab^2=\pm5$ and $x=3$.

\problem{8}
\hint{1} Look for a one-parameter family: take $m$ linear in $n$ so that $m-n$ divides both factors $m^2+n$ and $n^2+m$ in a controlled way. Small solutions found by computer, such as $(m,n)=(3,1),(5,2),(7,3)$, suggest the pattern.
\hint{2} With $m=2n+1$: $m^2+n=(4n+1)(n+1)$ and $n^2+m=(n+1)^2$, while $(m-n)^3=(n+1)^3$. So $a=4n+1$ is friendly for every $n\ge1$. This gives $a=5,9,\dots,2009$, which is $502$ values.
\hint{3} Now $a=2$. A solution needs $m>n$. Expand: the equation becomes $m^3-(n^2+6n)m^2+(6n^2-n)m-3n^3=0$. Put $g=\gcd(m,n)$, $m=gx$, $n=gy$ with $\gcd(x,y)=1$, $x>y$, and divide by $g^2$: $xy=gK$ with $K=x^3-gx^2y^2-6x^2y+6xy^2-3y^3$ (so $K>0$).
\hint{4} From the equation, $y\mid gx^3$, hence $y\mid g$; write $g=ys$. Then $x=sK$, and reducing $K$ mod $y$ gives $K\equiv x^3$, so $\gcd(K,y)=1$; also $x\mid3g=3ys$ gives $K\mid3y$. Hence $K\in\{1,3\}$, and $K=x^3\bigl(1-\frac{y^3}{K}\bigr)-6xy(x-y)-3y^3$.
\hint{5} If $K=1$, or $K=3$ and $y\ge2$, the right-hand side is negative. If $K=3$, $y=1$, you get $x^3-9x^2+9x-9=0$, which has no integer root (it would divide $9$). So $a=2$ is not friendly.

\problem{9}
\hint{1} Check the formula by listing representations: $r_4(n)=8,24,32,24,48,96,64,24,104,144$ for $n=1,\dots,10$, and continue up to $20$. The theta route uses Jacobi's identity $\theta(q)^4=1+8\sum_{4\nmid m}\frac{mq^m}{1-q^m}$, whose coefficient of $q^n$ is exactly $8\sum_{d\mid n,\,4\nmid d}d$. Proving that identity is the long part.
\hint{2} Four-square theorem: Euler's four-square identity shows that the representable numbers are closed under multiplication, so it suffices to treat primes $p$. By pigeonhole on the sets $\{x^2\}$ and $\{-1-y^2\}$ mod $p$ (each has $\frac{p+1}2$ elements), $x^2+y^2+1\equiv0$. Hence $mp$ is a sum of four squares for some $1\le m<p$.
\hint{3} Descent: let $m>1$ be minimal with $mp=\sum x_i^2$. Choose $y_i\equiv x_i\pmod m$ with $|y_i|\le m/2$; then $\sum y_i^2=mr$ with $r<m$ (and $r>0$). Multiply via Euler's identity: $m^2rp$ is a sum of four squares each divisible by $m$, so $rp$ is a sum of four squares, contradicting minimality. (If $m$ is even, you can first halve $m$ by pairing the $x_i$ by parity.)

\problem{10}
\hint{1} Cubes modulo $9$ are $0,\pm1$. So a sum of three cubes is never $\equiv\pm4\pmod9$, and this already gives infinitely many non-representable integers.
\hint{2} $33\equiv42\equiv6\pmod9$, and $6\equiv(-1)\cdot3$ is a sum of three cubic residues. Moreover, $9$ is the only modulus giving an obstruction (cubes cover enough residues modulo every other prime power), so congruences cannot decide these cases. In fact both were found to be sums of three cubes (Booker 2019; Booker--Sutherland 2019) after enormous computer searches.
\end{document}
